\documentclass[10pt,a4paper]{amsart}
\usepackage[margin=19mm]{geometry}
\usepackage{amsmath,amssymb,mathtools}
\usepackage[scr=rsfs,cal=euler]{mathalfa}
\usepackage{xfrac}
\usepackage[unicode,hidelinks,bookmarks=false]{hyperref}
\newcommand{\ie}{i.e.\ }
\newcommand{\cf}{cf.\ }
\newcommand{\resp}{respectively\ }
\newcommand{\Subsection}{\subsection}
\providecommand{\nobreakdash}{\nobreak\hbox{-}}
\setlength{\emergencystretch}{2em}
\multlinegap0pt
\usepackage{amssymb}
\newtheorem{lemma}{Lemma}
\title{Low-regularity well-posedness for the ZK Equation on a half-strip}
\author{E Avelino, G Doronin}
\date{Source: arXiv:2605.23092v1 (CC BY 4.0)}
\begin{document}
\maketitle

\noindent\emph{Source and licence.} Low-regularity well-posedness for the ZK Equation on a half-strip. Version arXiv:2605.23092v1 (CC BY 4.0), distributed by arXiv under the Creative Commons Attribution 4.0 International licence, http://creativecommons.org/licenses/by/4.0/, as recorded in the arXiv licence metadata for this exact version and shown on the abstract page https://arxiv.org/abs/2605.23092v1. Full licence notice: \texttt{licenses/arxiv-2605.23092-CC-BY-4.0.txt}.\par\smallskip

\noindent\textit{Editorial scope.} Counted unit: the article's lemma bilinear (source0.tex lines 430--435). The article's complete original proof of it is delivered with it (source0.tex lines 436--556, one \texttt{proof} environment of 121 lines). No mathematics is written, repaired or re-derived: the statement and the proof are the article's own lines, in the article's own notation, and no retained line is substituted or re-flowed.  No further source-local context is needed: every cross-reference of every retained line resolves inside the two delivered spans.  Nothing was omitted from any retained span, so the package carries no omission record.  No retained line contains a citation, so the wrapper carries no bibliography and no bibliography entry.  No retained line contains a figure environment, an image-inclusion command or drawing code, so no image is delivered and the package declares no asset dependency.  Editorial formatting only, in addition: an AMS \texttt{amsart} preamble that restates the article's own declarations and macros used by the retained lines, namely the environments \texttt{lemma} on separate counters, together with the standard packages those lines need.  The retained environments print in this document as Lemma 1 in order of appearance, whereas the article prints the same items with the numbering of its own full document.  The complete native source of the article as distributed in the arXiv e-print for this exact version is bundled unchanged as \texttt{sources/201182/source0.tex}.
\begin{lemma}
For $v \in Z^{b} = X^{0,b} \cap L^{2}_{t} H^{1}_{x,y}$, we have, for suitable $3/8<b<1/2$ and $1/2<\alpha<2/3$, the bilinear estimate
\begin{equation}\label{bilinear}
\|\partial_{x}(v^{2})\|_{Y^{0,b}}\leq C\|v\|_{Z^{b}}^{2}.
\end{equation}
\end{lemma}

\begin{proof}
Let $f = \partial_{x}(v^2)$. Thus, we have $\widehat{f}_{k}(\xi,\tau) = i\xi \mathcal{F}(v^{2})_{k}(\xi,\tau)$. To prove \eqref{bilinear}, we must estimate the three components of the $Y^{0,b}$ norm. To this end, we split the domain of integration into two regions: $|\omega_{k}(\xi)| < 1$ and $|\omega_{k}(\xi)| \ge 1$.

\textbf{Case A.} $|\omega_{k}(\xi)| < 1$.

In this region, the structure of $\omega_{k}(\xi)$ ensures that the
domain is bounded in the spatial frequency variable; that is, there
exists a constant $M_0 > 0$ such that $|\xi| \le M_0$. Consider the
second term of the $Y^{0,b}$ norm given by
\begin{equation}I_{low} = \sum_{k=1}^{\infty} \iint_{|\omega_{k}| < 1} \frac{|\xi|^{2} |\mathcal{F}(v^{2})_{k}(\xi,\tau)|^{2}}{\langle \tau \rangle^{2(1-\alpha)}} d\xi d\tau.\end{equation}
Since $\langle \tau \rangle^{-2(1-\alpha)} \le 1$ for all $\tau \in
\mathbb{R}$, and using the bound on the spatial frequency, we
estimate $I_{low}$ by
\begin{equation}
    I_{low} \le M_0^2 \sum_{k=1}^{\infty} \iint_{\mathbb{R}^2} |\mathcal{F}(v^2)_k(\xi,\tau)|^2 d\xi d\tau.
\end{equation}
By Plancherel's theorem, we obtain that
\begin{equation}
    I_{low} \le M_{0}^{2} \|v^{2}\|_{L^{2}_{x,y,t}}^{2} = M_{0}^{2} \|v\|_{L^{4}_{x,y,t}}^{4}.
\end{equation}
By the Strichartz estimate for the ZK equation, the space $X^{0,b}$
embeds continuously into $L^{4}_{x,y,t}$ provided that $b > 3/8$.
Since $b \in (3/8, 1/2)$, there exists a constant $C_{S} > 0$ such
that
\begin{equation}
    I_{low} \le M_{0}^{2} C_{S}^{4} \|v\|_{X^{0,b}}^{4} \le C_{1} \|v\|_{Z^{b}}^{4}.
\end{equation}

\textbf{Case B.} $|\omega_{k}(\xi)| \ge 1$.

The contribution of this region to the $Y^{0,b}$ norm decomposes
into two terms, $I_{high,1}$ and $I_{high,2}$, to be analyzed
separately via duality arguments.

\textbf{Case B1.} Estimate for $I_{high,1}$.

Let the corresponding term of the $Y^{0,b}$ norm be given by
\begin{equation}
    I_{high,1} = \sum_{k=1}^{\infty} \iint_{|\omega_{k}(\xi)| \ge 1} \frac{|\xi|^{2}
    |\mathcal{F}(v^{2})_{k}(\xi,\tau)|^{2}}{\langle \tau - \omega_{k}(\xi)\rangle^{2b}} d\xi d\tau.
\end{equation}
By duality in $L^2$, we introduce a non-negative test function $g_{k}(\xi,\tau)$ with $\|g\|_{L^2} \le 1$.
We must evaluate the linear functional
\begin{equation}\label{6.6}
    J_{1} = \sum_{k=1}^{\infty} \iint_{|\omega_{k}(\xi)|\geq 1} \frac{|\xi| |\mathcal{F}(v^{2})_k(\xi,\tau)|}{\langle \tau
    - \omega_{k}(\xi)\rangle^{b}} g_{k}(\xi,\tau) d\xi d\tau.
\end{equation}
Defining the positive function $\widehat{V}_{j} = |\widehat{v}_{j}|$
and expanding the product into a convolution, we bound $J_{1}$ by
its extended form over $\mathbb{R}^{3}$. By the symmetry of the
convolution, we integrate over the region where $|\xi_{1}| \ge
|\xi_{2}|$, which implies $|\xi| \le |\xi_{1}| + |\xi_{2}| \le
2|\xi_{1}|$. Hence, \eqref{6.6} becomes
\begin{equation}
    J_{1} \leq C\sum_{k,k_{1}} \iiint_{\mathbb{R}^{3}} \left[ \frac{g_{k}(\xi,\tau)}{\langle \tau -
    \omega_{k}(\xi)\rangle^{b}} \right] \cdot \Big( |\xi_{1}| \widehat{V}_{k_{1}}(\xi_{1},\tau_{1}) \Big) \cdot \widehat{V}_{k_{2}}(\xi_{2},\tau_{2}) d\Gamma,
\end{equation}
where $d\Gamma$ denotes the measure under the constraint $\xi =
\xi_{1}+\xi_{2}$ and $\tau = \tau_{1}+\tau_{2}$. Applying
Plancherel's theorem and the generalized Hölder inequality with
indices $1/4 + 1/2 + 1/4 = 1$, we have
\begin{equation}
J_{1} \leq C\left\|\mathcal{F}^{-1}\left(\frac{g_{k}}{\langle \tau - \omega_{k}(\xi)
\rangle^{b}}\right)\right\|_{L^{4}} \cdot \left\|\mathcal{F}^{-1}\left(|\xi_1|
\widehat{V}_{k_{1}}\right)\right\|_{L^{2}} \cdot \left\|\mathcal{F}^{-1}\left(\widehat{V}_{k_{2}}\right)\right\|_{L^{4}}.
\end{equation}
Hence,
\begin{equation}
\left\|\mathcal{F}^{-1}\left(|\xi_1| \widehat{V}_{k_{1}}\right)\right\|_{L^{2}}= \|\partial_{x} v\|_{L^{2}} \leq \|v\|_{L^{2}_{t} H^{1}_{x,y}} \leq \|v\|_{Z^{b}},
\end{equation}
and
\begin{equation}
\left\|\mathcal{F}^{-1}\left(\widehat{V}_{k_{2}}\right)\right\|_{L^{4}} \leq C_{S} \|v\|_{X^{0,b}} \leq C_{S} \|v\|_{Z^{b}}.
\end{equation}
Since $\|g\|_{L^{2}} \le 1$, we have that $\mathcal{F}^{-1}\left(\frac{g_{k}}{\langle \tau - \omega_{k}(\xi) \rangle^{b}}\right) \in X^{0,b}$. Since $b > 3/8$, we apply the Strichartz embedding to obtain
\begin{equation}
\left\|\mathcal{F}^{-1}\left(\frac{g_{k}}{\langle \tau - \omega_{k}(\xi) \rangle^{b}}\right)\right\|_{L^4} \le C_{S}\left\|\mathcal{F}^{-1}\left(\frac{g_{k}}{\langle \tau - \omega_{k}(\xi) \rangle^{b}}\right)\right\|_{X^{0,b}} \le 1.\end{equation}
Therefore, $J_{1} \le C_{h,1}\|v\|_{Z^{b}}^{2}$, which implies $I_{high,1} \le C_{2}\|v\|_{Z^{b}}^{4}$.

\textbf{Case B2.} Estimate for $I_{high,2}$.

Let the corresponding term of the $Y^{0,b}$ norm be given by
\begin{equation}
I_{high,2} = \sum_{k=1}^{\infty} \int_{|\omega_{k}| \ge 1} \left( \int_{|\omega_{k}| \ge 1} \frac{|\xi| |\mathcal{F}(v^{2})_{k}(\xi,\tau)|}{\langle \tau - \omega_{k}(\xi) \rangle} d\tau \right)^{2} d\xi.
\end{equation}
By duality in the mixed norm $L^{2}_{\xi, k}$, we take test
functions $h_{k}(\xi)$ depending only on the spatial frequency, such
that $\|h\|_{L^{2}} \le 1$. Arguing via symmetry breaking, where
$|\xi| \le 2|\xi_{1}|$, and using Hölder's inequality as in case
B1, we get
\begin{equation}
I_{high,2}^{1/2} \le C\left\|\mathcal{F}^{-1} \left( \frac{h_{k}(\xi)}{\langle \tau - \omega_{k}(\xi) \rangle} \right)\right\|_{L^{4}_{x,y,t}} \cdot \left\|\mathcal{F}^{-1}\left(|\xi_1| \widehat{V}_{k_{1}}\right)\right\|_{L^{2}_{x,y,t}} \cdot \left\|\mathcal{F}^{-1}\left(\widehat{V}_{k_{2}}\right)\right\|_{L^{4}_{x,y,t}}.
\end{equation}
It remains to show that $\mathcal{F}^{-1} \left(
\frac{h_{k}(\xi)}{\langle \tau - \omega_{k}(\xi) \rangle} \right)
\in L^{4}_{x,y,t}$.









Fix $s$ such that $3/8 < s < 1/2$. The norm of $\mathcal{F}^{-1}
\left( \frac{h_{k}(\xi)}{\langle \tau - \omega_{k}(\xi) \rangle}
\right)$ in the space $X^{0,s}$ is given by
\begin{equation}
\left\|\mathcal{F}^{-1} \left( \frac{h_{k}(\xi)}{\langle \tau - \omega_{k}(\xi) \rangle} \right)\right\|_{X^{0,s}}^{2} = \sum_{k=1}^{\infty} \int_{\mathbb{R}} |h_{k}(\xi)|^{2} \left( \int_{\mathbb{R}} \frac{1}{\langle \tau - \omega_{k}(\xi) \rangle^{2(1 -s)}} d\tau \right) d\xi.
\end{equation}
Since $s < 1/2$, we have $2(1 - s) > 1$. Thus, the inner integral with respect to $\tau$ converges to a constant $C_{s} > 0$. It follows that
\begin{equation}
\left\|\mathcal{F}^{-1} \left( \frac{h_{k}(\xi)}{\langle \tau - \omega_{k}(\xi) \rangle} \right)\right\|_{X^{0,s}}^{2} \le C_{s} \sum_{k=1}^{\infty} \int_{\mathbb{R}} |h_{k}(\xi)|^2 d\xi = C_{s} \|h\|_{L^{2}_{\xi, k}}^{2} \le C_{s}.
\end{equation}
Since $\mathcal{F}^{-1} \left( \frac{h_{k}(\xi)}{\langle \tau - \omega_{k}(\xi) \rangle} \right) \in X^{0,s}$ with $s > 3/8$, the Strichartz embedding guarantees that
\begin{equation}\left\|\mathcal{F}^{-1} \left( \frac{h_{k}(\xi)}{\langle \tau - \omega_{k}(\xi) \rangle} \right)\right\|_{L^{4}} \le C_{s}.\end{equation}
Since the factors $\mathcal{F}^{-1}\left(|\xi_{1}| \widehat{V}_{k_{1}}\right)$ and $\mathcal{F}^{-1}\left(\widehat{V}_{k_{2}}\right)$ are bounded by $\|v\|_{Z^{b}}$ and $C_{S}\|v\|_{Z^{b}}$ respectively, we conclude that $I_{high,2}^{1/2} \le C_{h,2}\|v\|_{Z^{b}}^{2}$, that is, $I_{high,2} \le C_{3}\|v\|_{Z^{b}}^{4}$.

Combining the estimates for $I_{low}$, $I_{high,1}$, and $I_{high,2}$, the result follows.
\end{proof}

\end{document}
